% GENERATED FILE. Edit canonical lessons, then run npm run generate:books.
% canonical-source-hash: fnv1a64:df5853950d683082
% canonical-lessons: MW-C289-masin, MW-C289-chalavno, MW-C289-ghisno, MW-C289-rando, MW-C289-chitarno

\chapter{Tools and Machines}
\label{ch:mw-tools-and-machines}
\hlchaptermodality{289}

\begin{chapteropening}
\textbf{By the end of this chapter:} \emph{I can say what tools and machines I can use.}

Complete the last lesson of chapter 289: I can say what tools and machines I can use.
\end{chapteropening}

\section[\texorpdfstring{\mw{मसीन} (masīn)}{मसीन (masīn)}]{\mw{मसीन} (masīn) — a machine}

\label{lesson:MW-C289-masin}



\begin{quote}
Before the new one: say the Marwadi for a trade, then the Marwadi for a labourer.
\end{quote}

\subsection*{You'll want to know: \mw{मसीन}}
\textbf{\mw{मसीन}} — \emph{masīn} — \textquotedblleft{}a machine\textquotedblright{}.

\begin{quote}
\mw{मसीन} \mw{नुवी} \mw{है।}
\end{quote}

\begin{quote}
\emph{masīn nuvī hai.}
\end{quote}

\begin{quote}
The machine is new.
\end{quote}

A feminine noun.

\begin{cousinweb}
\textbf{\mw{मसीन}} (\emph{masīn}) is borrowed from English \emph{machine}. It is feminine: \textbf{\mw{मसीन} \mw{नुवी} \mw{है}}.
\end{cousinweb}

\begin{grammarlens}[title={What you've built}]
One more word for this chapter.
\end{grammarlens}

\subsection*{Your turn}
\begin{itemize}
\raggedright
  \item \emph{Say it:} \emph{masīn}
  \item \emph{Say it:} \emph{masīn}, once more
  \item \emph{Say it:} \emph{masīn nuvī hai.} — The machine is new.
  \item \emph{Recall:} say the Marwadi for a trade, then the Marwadi for a labourer, then say \emph{masīn} again
\end{itemize}

\begin{tcolorbox}[breakable,colback=teal!4,colframe=teal!35!black,title={Before you move on}]
What does \mw{मसीन} mean? (\textquotedblleft{}A machine\textquotedblright{}.) Say it once more.
\end{tcolorbox}
\section[\texorpdfstring{\mw{चलावणो} (chalāvṇo)}{चलावणो (chalāvṇo)}]{\mw{चलावणो} (chalāvṇo) — to drive; to run (a machine)}

\label{lesson:MW-C289-chalavno}



\begin{quote}
Before the new one: say the Marwadi for a labourer, then the Marwadi for a machine.
\end{quote}

\subsection*{You'll want to know: \mw{चलावणो}}
\textbf{\mw{चलावणो}} — \emph{chalāvṇo} — \textquotedblleft{}to drive; to run (a machine)\textquotedblright{}.

\begin{quote}
\mw{भाई} \mw{मसीन} \mw{चलावै} \mw{है।}
\end{quote}

\begin{quote}
\emph{bhāī masīn chalāvai hai.}
\end{quote}

\begin{quote}
My brother runs the machine.
\end{quote}

Past (I …): \mw{म्हैं} \mw{गाड़ी} \mw{चलाई}, \emph{mhaĩ gāṛī chalāī}.

\begin{cousinweb}
\textbf{\mw{चलावणो}} (\emph{chalāvṇo}) is built on \textbf{\mw{चालणो}}, to walk, as its causative: \textquotedblleft{}to make go\textquotedblright{}. CDIAL lists Old Marwari \emph{calaï}, \textquotedblleft{}goes\textquotedblright{}, under Sanskrit \emph{calati} (CDIAL 4715).
\end{cousinweb}

\begin{grammarlens}[title={What you've built}]
One more word for this chapter.
\end{grammarlens}

\subsection*{Your turn}
\begin{itemize}
\raggedright
  \item \emph{Say it:} \emph{chalāvṇo}
  \item \emph{Say it:} \emph{chalāvṇo}, once more
  \item \emph{Say it:} \emph{bhāī masīn chalāvai hai.} — My brother runs the machine.
  \item \emph{Recall:} say the Marwadi for a labourer, then the Marwadi for a machine, then say \emph{chalāvṇo} again
\end{itemize}

\begin{tcolorbox}[breakable,colback=teal!4,colframe=teal!35!black,title={Before you move on}]
What does \mw{चलावणो} mean? (\textquotedblleft{}To drive; to run (a machine)\textquotedblright{}.) Say it once more.
\end{tcolorbox}
\section[\texorpdfstring{\mw{घिसणो} (ghisṇo)}{घिसणो (ghisṇo)}]{\mw{घिसणो} (ghisṇo) — to rub smooth; to rub}

\label{lesson:MW-C289-ghisno}



\begin{quote}
Before the new one: say the Marwadi for a machine, then the Marwadi for to drive.
\end{quote}

\subsection*{You'll want to know: \mw{घिसणो}}
\textbf{\mw{घिसणो}} — \emph{ghisṇo} — \textquotedblleft{}to rub smooth; to rub\textquotedblright{}.

\begin{quote}
\mw{काको} \mw{लकड़ी} \mw{घिसै} \mw{है।}
\end{quote}

\begin{quote}
\emph{kāko lakṛī ghisai hai.}
\end{quote}

\begin{quote}
My uncle rubs the wood smooth.
\end{quote}

Past (I …): \mw{म्हैं} \mw{लकड़ी} \mw{घिसी}, \emph{mhaĩ lakṛī ghisī}.

\begin{cousinweb}
\textbf{\mw{घिसणो}} (\emph{ghisṇo}) is inherited from Sanskrit \emph{gharṣati}, \textquotedblleft{}rubs\textquotedblright{}; CDIAL gives no Marwari form in this entry, but lists Hindi \emph{ghisnā}, \textquotedblleft{}to rub\textquotedblright{}, and Gujarati \emph{ghasvũ}, \textquotedblleft{}to rub, polish\textquotedblright{} (CDIAL 4450).
\end{cousinweb}

\begin{grammarlens}[title={What you've built}]
One more word for this chapter.
\end{grammarlens}

\subsection*{Your turn}
\begin{itemize}
\raggedright
  \item \emph{Say it:} \emph{ghisṇo}
  \item \emph{Say it:} \emph{ghisṇo}, once more
  \item \emph{Say it:} \emph{kāko lakṛī ghisai hai.} — My uncle rubs the wood smooth.
  \item \emph{Recall:} say the Marwadi for a machine, then the Marwadi for to drive, then say \emph{ghisṇo} again
\end{itemize}

\begin{tcolorbox}[breakable,colback=teal!4,colframe=teal!35!black,title={Before you move on}]
What does \mw{घिसणो} mean? (\textquotedblleft{}To rub smooth; to rub\textquotedblright{}.) Say it once more.
\end{tcolorbox}
\section[\texorpdfstring{\mw{रंदो} (rando)}{रंदो (rando)}]{\mw{रंदो} (rando) — a carpenter's plane}

\label{lesson:MW-C289-rando}



\begin{quote}
Before the new one: say the Marwadi for to drive, then the Marwadi for to rub smooth.
\end{quote}

\subsection*{You'll want to know: \mw{रंदो}}
\textbf{\mw{रंदो}} — \emph{rando} — \textquotedblleft{}a carpenter's plane\textquotedblright{}.

\begin{quote}
\mw{रंदो} \mw{भारी} \mw{है।}
\end{quote}

\begin{quote}
\emph{rando bhārī hai.}
\end{quote}

\begin{quote}
The plane is heavy.
\end{quote}

A masculine noun.

\begin{cousinweb}
\textbf{\mw{रंदो}} (\emph{rando}) is usually traced to Persian \emph{randa}, \textquotedblleft{}a carpenter's plane\textquotedblright{}.
\end{cousinweb}

\begin{grammarlens}[title={What you've built}]
One more word for this chapter.
\end{grammarlens}

\subsection*{Your turn}
\begin{itemize}
\raggedright
  \item \emph{Say it:} \emph{rando}
  \item \emph{Say it:} \emph{rando}, once more
  \item \emph{Say it:} \emph{rando bhārī hai.} — The plane is heavy.
  \item \emph{Recall:} say the Marwadi for to drive, then the Marwadi for to rub smooth, then say \emph{rando} again
\end{itemize}

\begin{tcolorbox}[breakable,colback=teal!4,colframe=teal!35!black,title={Before you move on}]
What does \mw{रंदो} mean? (\textquotedblleft{}A carpenter's plane\textquotedblright{}.) Say it once more.
\end{tcolorbox}
\section[\texorpdfstring{\mw{चीतरणो} (chītarṇo)}{चीतरणो (chītarṇo)}]{\mw{चीतरणो} (chītarṇo) — to paint a picture}

\label{lesson:MW-C289-chitarno}



\begin{quote}
Before the new one: say the Marwadi for to rub smooth, then the Marwadi for a carpenter's plane.
\end{quote}

\subsection*{You'll want to know: \mw{चीतरणो}}
\textbf{\mw{चीतरणो}} — \emph{chītarṇo} — \textquotedblleft{}to paint a picture\textquotedblright{}.

\begin{quote}
\mw{छोरी} \mw{भींत} \mw{माथै} \mw{मोर} \mw{चीतरै} \mw{है।}
\end{quote}

\begin{quote}
\emph{chhorī bhī̃t māthai mor chītrai hai.}
\end{quote}

\begin{quote}
The girl paints a peacock on the wall.
\end{quote}

Past (I …): \mw{म्हैं} \mw{मोर} \mw{चीतर्यो}, \emph{mhaĩ mor chītaryo}.

\begin{cousinweb}
\textbf{\mw{चीतरणो}} (\emph{chītarṇo}) is inherited from Sanskrit \emph{citrakāra}, \textquotedblleft{}a painter\textquotedblright{}; CDIAL gives no Marwari form in this entry, but lists Hindi \emph{citārnā}, \textquotedblleft{}to paint\textquotedblright{} (CDIAL 4807).
\end{cousinweb}

\begin{grammarlens}[title={What you've built}]
One more word for this chapter.
\end{grammarlens}

\subsection*{Your turn}
\begin{itemize}
\raggedright
  \item \emph{Say it:} \emph{chītarṇo}
  \item \emph{Say it:} \emph{chītarṇo}, once more
  \item \emph{Say it:} \emph{chhorī bhī̃t māthai mor chītrai hai.} — The girl paints a peacock on the wall.
  \item \emph{Recall:} say the Marwadi for to rub smooth, then the Marwadi for a carpenter's plane, then say \emph{chītarṇo} again
\end{itemize}

\begin{tcolorbox}[breakable,colback=teal!4,colframe=teal!35!black,title={Before you move on}]
What does \mw{चीतरणो} mean? (\textquotedblleft{}To paint a picture\textquotedblright{}.) Say it once more.
\end{tcolorbox}
